%==============================================================================
%  FIGURE SPECIFICATIONS FOR ILLUSTRATOR
%  Optimization for AI: From Gradient Descent to Modern Optimizers
%  Book 2 - The AI Engineer's Library
%
%  This document provides detailed specifications for all figures needed
%  in Chapters 9 (Convex Optimization) and 15 (Convergence Analysis).
%==============================================================================

\documentclass[12pt]{article}
\usepackage[margin=1in]{geometry}
\usepackage{booktabs}
\usepackage{longtable}
\usepackage{xcolor}
\usepackage{hyperref}
\usepackage{enumitem}

% Color palette for the book
\definecolor{primaryblue}{RGB}{31, 119, 180}
\definecolor{secondaryorange}{RGB}{255, 127, 14}
\definecolor{successgreen}{RGB}{44, 160, 44}
\definecolor{dangerred}{RGB}{214, 39, 40}
\definecolor{warningyellow}{RGB}{255, 193, 7}
\definecolor{accentpurple}{RGB}{148, 103, 189}
\definecolor{neutralgray}{RGB}{127, 127, 127}
\definecolor{lightblue}{RGB}{174, 199, 232}
\definecolor{lightorange}{RGB}{255, 187, 120}

\title{Figure Specifications for Illustrator\\[0.5em]
\large Optimization for AI: From Gradient Descent to Modern Optimizers\\
Book 2 -- The AI Engineer's Library}
\author{Dr. Vijay Raghavan}
\date{April 2026}

\begin{document}

\maketitle

\tableofcontents
\newpage

%==============================================================================
\section{General Guidelines}
%==============================================================================

\subsection{Color Palette}
All figures should use the following consistent color palette:
\begin{itemize}
    \item \textbf{Primary Blue} (\#1F77B4): Main elements, positive examples, optimal paths
    \item \textbf{Secondary Orange} (\#FF7F0E): Contrasting elements, alternative methods
    \item \textbf{Success Green} (\#2CA02C): Good outcomes, convex regions, convergence
    \item \textbf{Danger Red} (\#D62728): Bad outcomes, non-convex regions, divergence
    \item \textbf{Accent Purple} (\#9467BD): Tertiary elements, annotations
    \item \textbf{Neutral Gray} (\#7F7F7F): Grid lines, axes, secondary annotations
    \item \textbf{Light Blue} (\#AEC7E8): Shaded regions, fills
    \item \textbf{Light Orange} (\#FFBB78): Secondary fills
\end{itemize}

\subsection{Typography}
\begin{itemize}
    \item \textbf{Axis labels}: 10pt sans-serif (e.g., Helvetica, Arial)
    \item \textbf{Tick labels}: 8pt sans-serif
    \item \textbf{Annotations}: 9pt sans-serif
    \item \textbf{Mathematical notation}: Use proper LaTeX-style symbols ($\theta$, $\nabla$, etc.)
\end{itemize}

\subsection{Dimensions}
\begin{itemize}
    \item \textbf{Full-width figures}: 6.0 inches (15.24 cm) wide
    \item \textbf{Side-by-side figures}: 2.8 inches (7.11 cm) each
    \item \textbf{Resolution}: 300 DPI minimum for print
    \item \textbf{Format}: PDF (vector) preferred, PNG for raster elements
\end{itemize}

\subsection{Accessibility}
\begin{itemize}
    \item Use distinct line styles (solid, dashed, dotted) in addition to colors
    \item Ensure sufficient contrast (4.5:1 minimum)
    \item Include descriptive alt-text for each figure (provided below)
\end{itemize}

\newpage

%==============================================================================
\section{Chapter 9: Convex Optimization (8 Figures)}
%==============================================================================

%------------------------------------------------------------------------------
\subsection{Figure 9.1: Convex vs Non-Convex Sets}
%------------------------------------------------------------------------------

\begin{description}[style=nextline, leftmargin=0pt]
    \item[Figure ID] \texttt{fig:convex-nonconvex-sets}
    \item[Title] Convex vs Non-Convex Sets: The Line Segment Test
    \item[Type] 2D geometric illustration (side-by-side comparison)
    \item[Size] 6.0 $\times$ 3.0 inches (full width, two panels)

    \item[Visual Description]
    Two panels side by side:

    \textbf{Left Panel -- ``Convex Set'':}
    \begin{itemize}
        \item Draw a smooth, rounded convex shape (ellipse or rounded polygon)
        \item Fill with Light Blue (\#AEC7E8) at 40\% opacity
        \item Outline in Primary Blue (\#1F77B4), 2pt line weight
        \item Mark two arbitrary points $\mathbf{x}$ and $\mathbf{y}$ inside the shape as filled circles (Primary Blue)
        \item Draw a straight line segment connecting $\mathbf{x}$ and $\mathbf{y}$ (Success Green, 2pt, dashed)
        \item The entire line segment should be visibly inside the shape
        \item Add a checkmark symbol or ``$\checkmark$'' in Success Green
    \end{itemize}

    \textbf{Right Panel -- ``Non-Convex Set'':}
    \begin{itemize}
        \item Draw a crescent moon or star shape (clearly non-convex with indentation)
        \item Fill with Light Orange (\#FFBB78) at 40\% opacity
        \item Outline in Secondary Orange (\#FF7F0E), 2pt line weight
        \item Mark two points $\mathbf{x}$ and $\mathbf{y}$ on opposite sides of the indentation
        \item Draw a straight line segment connecting them (Danger Red, 2pt, dashed)
        \item The line segment should clearly pass outside the shape through the indentation
        \item Highlight the portion outside the set with a different pattern or thicker line
        \item Add an ``$\times$'' symbol in Danger Red where the line exits the set
    \end{itemize}

    \item[Required Labels]
    \begin{itemize}
        \item Panel titles: ``Convex Set'' (left), ``Non-Convex Set'' (right)
        \item Point labels: $\mathbf{x}$, $\mathbf{y}$ (both panels)
        \item Left panel annotation: ``Line segment stays inside''
        \item Right panel annotation: ``Line segment exits the set''
    \end{itemize}

    \item[Caption Text]
    ``The line segment test for convexity. A set is convex if and only if the line segment connecting any two points in the set lies entirely within the set. Left: An ellipse is convex---every possible line segment stays inside. Right: A crescent is non-convex---some line segments pass outside the set.''

    \item[Alt-Text for Accessibility]
    ``Side-by-side comparison showing convex and non-convex sets. Left panel shows a blue ellipse with two points connected by a green dashed line that stays entirely inside the shape. Right panel shows an orange crescent shape with two points connected by a red dashed line that passes outside the shape through the indentation.''
\end{description}

\newpage

%------------------------------------------------------------------------------
\subsection{Figure 9.2: Convex vs Concave Functions}
%------------------------------------------------------------------------------

\begin{description}[style=nextline, leftmargin=0pt]
    \item[Figure ID] \texttt{fig:convex-concave-3d}
    \item[Title] Convex vs Concave Functions: 3D Bowl Comparison
    \item[Type] 3D surface plots (side-by-side)
    \item[Size] 6.0 $\times$ 3.5 inches (full width, two panels)

    \item[Visual Description]
    Two 3D surface plots side by side:

    \textbf{Left Panel -- ``Convex Function (Bowl)'':}
    \begin{itemize}
        \item Plot $f(x, y) = x^2 + y^2$ as a 3D surface
        \item Domain: $x, y \in [-2, 2]$
        \item Color gradient from Primary Blue (low) to Light Blue (high)
        \item Show the surface as a bowl opening upward
        \item Mark the minimum point at the origin with a red dot
        \item Draw two points on the surface with a chord (line segment) connecting them
        \item The chord should be visibly ABOVE the surface (illustrating $f(\lambda x + (1-\lambda)y) \leq \lambda f(x) + (1-\lambda)f(y)$)
        \item Use subtle grid lines on the surface
    \end{itemize}

    \textbf{Right Panel -- ``Concave Function (Dome)'':}
    \begin{itemize}
        \item Plot $f(x, y) = -x^2 - y^2$ (or $4 - x^2 - y^2$ for better visualization)
        \item Domain: $x, y \in [-2, 2]$
        \item Color gradient from Secondary Orange (low) to Light Orange (high)
        \item Show the surface as a dome opening downward (inverted bowl)
        \item Mark the maximum point at the origin with a green dot
        \item Draw two points on the surface with a chord connecting them
        \item The chord should be visibly BELOW the surface
        \item Use subtle grid lines on the surface
    \end{itemize}

    \item[Required Labels]
    \begin{itemize}
        \item Panel titles: ``Convex: $f(x,y) = x^2 + y^2$'' (left), ``Concave: $f(x,y) = -x^2 - y^2$'' (right)
        \item Axes: $x$, $y$, $f(x,y)$
        \item Left annotation: ``Chord lies above surface''
        \item Right annotation: ``Chord lies below surface''
        \item Left: ``Global minimum'' pointing to origin
        \item Right: ``Global maximum'' pointing to origin
    \end{itemize}

    \item[Caption Text]
    ``Convex and concave functions in 3D. Left: A convex function (paraboloid opening upward) where any chord connecting two points on the surface lies above or on the surface. The unique global minimum is at the bottom of the bowl. Right: A concave function (paraboloid opening downward) where chords lie below the surface, with a unique global maximum at the top.''

    \item[Alt-Text for Accessibility]
    ``Two 3D surface plots side by side. Left shows a blue bowl-shaped paraboloid opening upward with a line segment (chord) connecting two points on the surface, the chord floating above the surface. Right shows an orange dome-shaped paraboloid opening downward with a chord connecting two points, the chord positioned below the surface.''
\end{description}

\newpage

%------------------------------------------------------------------------------
\subsection{Figure 9.3: First-Order Condition for Convexity}
%------------------------------------------------------------------------------

\begin{description}[style=nextline, leftmargin=0pt]
    \item[Figure ID] \texttt{fig:first-order-condition}
    \item[Title] First-Order Condition: Tangent Plane as Global Underestimator
    \item[Type] 3D surface plot with tangent plane
    \item[Size] 5.0 $\times$ 4.0 inches

    \item[Visual Description]
    \begin{itemize}
        \item Plot a convex function surface: $f(x, y) = x^2 + 0.5y^2$ or similar
        \item Domain: $x, y \in [-2, 2]$
        \item Surface color: Primary Blue with transparency (70\% opacity)
        \item Choose a point $\mathbf{x}_0$ on the surface (e.g., at $(1, 0.5)$)
        \item Mark this point with a prominent red dot
        \item Draw the tangent plane at $\mathbf{x}_0$:
        \begin{itemize}
            \item Color: Success Green with 50\% transparency
            \item The plane should extend enough to show it lies entirely below/on the surface
            \item Edge the plane with a darker green line
        \end{itemize}
        \item Draw the gradient vector $\nabla f(\mathbf{x}_0)$ as an arrow in the tangent plane (Secondary Orange, thick arrow)
        \item Show that the tangent plane touches the surface only at $\mathbf{x}_0$
        \item Add visual gap annotations showing the tangent plane is below the surface elsewhere
    \end{itemize}

    \item[Required Labels]
    \begin{itemize}
        \item Point label: $\mathbf{x}_0$ at the tangent point
        \item Gradient arrow label: $\nabla f(\mathbf{x}_0)$
        \item Surface label: $f(\mathbf{x})$
        \item Tangent plane label: $f(\mathbf{x}_0) + \nabla f(\mathbf{x}_0)^\top (\mathbf{x} - \mathbf{x}_0)$
        \item Annotation: ``Tangent plane never exceeds surface''
        \item Axes: $x_1$, $x_2$, $f$
    \end{itemize}

    \item[Caption Text]
    ``First-order condition for convexity. For a convex function, the tangent plane at any point $\mathbf{x}_0$ lies entirely below (or on) the function surface. Mathematically: $f(\mathbf{y}) \geq f(\mathbf{x}_0) + \nabla f(\mathbf{x}_0)^\top (\mathbf{y} - \mathbf{x}_0)$ for all $\mathbf{y}$. This property means the linear approximation is always a global underestimator.''

    \item[Alt-Text for Accessibility]
    ``3D plot showing a blue bowl-shaped convex surface with a green tangent plane touching the surface at a single point marked with a red dot. The tangent plane extends beneath the entire surface, never crossing above it. An orange arrow shows the gradient direction at the tangent point.''
\end{description}

\newpage

%------------------------------------------------------------------------------
\subsection{Figure 9.4: Gradient Descent on Convex Function}
%------------------------------------------------------------------------------

\begin{description}[style=nextline, leftmargin=0pt]
    \item[Figure ID] \texttt{fig:gd-convex-contours}
    \item[Title] Gradient Descent on a Convex Function
    \item[Type] 2D contour plot with optimization trajectory
    \item[Size] 5.0 $\times$ 4.5 inches

    \item[Visual Description]
    \begin{itemize}
        \item Draw elliptical contour lines for $f(\theta_1, \theta_2) = \theta_1^2 + 4\theta_2^2$ (or similar)
        \item Domain: $\theta_1, \theta_2 \in [-3, 3]$
        \item Contour lines: Neutral Gray, varying line weights (thicker for lower values)
        \item Number of contour levels: 8--10 concentric ellipses
        \item Color-fill the contours with a blue gradient (darker at center/minimum)
        \item Mark the global minimum at origin with a star symbol (Success Green)
        \item Draw the gradient descent trajectory:
        \begin{itemize}
            \item Starting point: top-right corner, e.g., $(2.5, 1.5)$, marked with filled circle (Danger Red)
            \item Path: connected line segments showing 8--12 iterations
            \item Line style: Primary Blue, 2pt, with arrows showing direction
            \item Mark each iteration point with small circles
            \item Path should show characteristic perpendicular-to-contour steps
            \item Steps should get smaller near the minimum
        \end{itemize}
        \item Add gradient vectors at 2--3 points along the path (Secondary Orange arrows, pointing in descent direction)
    \end{itemize}

    \item[Required Labels]
    \begin{itemize}
        \item Axes: $\theta_1$, $\theta_2$
        \item Start point: ``Start: $\boldsymbol{\theta}_0$''
        \item End point: ``$\boldsymbol{\theta}^*$ (minimum)''
        \item Gradient arrows: $-\nabla f$
        \item Contour annotation: ``Level sets of $f(\boldsymbol{\theta})$''
        \item Iteration labels: $t=0$, $t=5$, $t=10$ (at selected points)
    \end{itemize}

    \item[Caption Text]
    ``Gradient descent trajectory on a convex quadratic function. The contour lines show level sets of the loss function, with the global minimum at the center (green star). Starting from the red point, gradient descent follows the negative gradient direction at each step, perpendicular to the contours. On convex functions, this path is guaranteed to reach the global minimum.''

    \item[Alt-Text for Accessibility]
    ``2D contour plot showing elliptical level curves of a convex function with a gradient descent path. The path starts from a red point in the upper right, follows a zigzag trajectory through the elliptical contours, and converges to a green star at the center representing the global minimum. Small orange arrows indicate the gradient direction at several points along the path.''
\end{description}

\newpage

%------------------------------------------------------------------------------
\subsection{Figure 9.5: SVM Decision Boundary}
%------------------------------------------------------------------------------

\begin{description}[style=nextline, leftmargin=0pt]
    \item[Figure ID] \texttt{fig:svm-decision-boundary}
    \item[Title] Support Vector Machine: Decision Boundary and Margin
    \item[Type] 2D scatter plot with separating hyperplane
    \item[Size] 5.0 $\times$ 4.0 inches

    \item[Visual Description]
    \begin{itemize}
        \item Create a 2D scatter plot with two linearly separable classes
        \item Class +1 (positive): 15--20 points as filled circles, Primary Blue
        \item Class -1 (negative): 15--20 points as filled triangles, Secondary Orange
        \item Ensure classes are clearly separated but not by too much
        \item Draw the optimal separating hyperplane (decision boundary):
        \begin{itemize}
            \item Solid black line, 2pt weight
            \item Position to maximize margin between classes
        \end{itemize}
        \item Draw margin boundaries (gutters):
        \begin{itemize}
            \item Two dashed lines parallel to the decision boundary
            \item Color: Neutral Gray, 1.5pt
            \item Positioned at $\mathbf{w}^\top \mathbf{x} + b = +1$ and $\mathbf{w}^\top \mathbf{x} + b = -1$
        \end{itemize}
        \item Shade the margin region between the gutters (light yellow or light gray, 20\% opacity)
        \item Identify and highlight support vectors:
        \begin{itemize}
            \item Circle the 3--5 points that lie exactly on the margin boundaries
            \item Use Danger Red circles around these points
            \item These should include points from both classes
        \end{itemize}
        \item Draw the margin width indicator:
        \begin{itemize}
            \item Double-headed arrow showing distance $\frac{2}{\|\mathbf{w}\|}$
            \item Perpendicular to the decision boundary
            \item Color: Success Green
        \end{itemize}
        \item Draw the weight vector $\mathbf{w}$ as an arrow perpendicular to the decision boundary (Accent Purple)
    \end{itemize}

    \item[Required Labels]
    \begin{itemize}
        \item Decision boundary: ``$\mathbf{w}^\top \mathbf{x} + b = 0$''
        \item Upper margin line: ``$\mathbf{w}^\top \mathbf{x} + b = +1$''
        \item Lower margin line: ``$\mathbf{w}^\top \mathbf{x} + b = -1$''
        \item Margin width: ``Margin $= \frac{2}{\|\mathbf{w}\|}$''
        \item Support vectors: ``Support vectors'' with arrow pointing to circled points
        \item Weight vector: $\mathbf{w}$
        \item Axes: $x_1$, $x_2$
        \item Legend: ``Class +1'' (blue circle), ``Class -1'' (orange triangle)
    \end{itemize}

    \item[Caption Text]
    ``Support Vector Machine decision boundary. The SVM finds the hyperplane that maximizes the margin (shaded region) between the two classes. The margin boundaries pass through the support vectors (circled in red)---the critical points that define the solution. All other points could be removed without changing the decision boundary.''

    \item[Alt-Text for Accessibility]
    ``2D scatter plot showing binary classification with SVM. Blue circles represent positive class points clustered in upper right. Orange triangles represent negative class points clustered in lower left. A solid black line separates the classes, with parallel dashed lines on each side forming the margin zone. Points touching the dashed margin lines are circled in red and labeled as support vectors. A green double-headed arrow indicates the margin width.''
\end{description}

\newpage

%------------------------------------------------------------------------------
\subsection{Figure 9.6: Hard Margin vs Soft Margin SVM}
%------------------------------------------------------------------------------

\begin{description}[style=nextline, leftmargin=0pt]
    \item[Figure ID] \texttt{fig:hard-soft-margin}
    \item[Title] Hard Margin vs Soft Margin SVM
    \item[Type] Side-by-side 2D scatter plots
    \item[Size] 6.0 $\times$ 3.0 inches (two panels)

    \item[Visual Description]
    Two panels showing the same dataset but with different SVM approaches:

    \textbf{Left Panel -- ``Hard Margin SVM'':}
    \begin{itemize}
        \item Show linearly separable data with one outlier/noise point
        \item The outlier from Class -1 intrudes into the Class +1 region
        \item Draw a decision boundary that perfectly separates all points
        \item This boundary will be skewed/suboptimal due to the outlier
        \item Margin will be very narrow
        \item Shade the narrow margin region
        \item Circle the support vectors
        \item Add annotation: ``Narrow margin due to outlier''
        \item Point to the outlier with label: ``Outlier''
    \end{itemize}

    \textbf{Right Panel -- ``Soft Margin SVM'':}
    \begin{itemize}
        \item Same data points as left panel
        \item Draw a more reasonable decision boundary that allows the outlier to be misclassified
        \item Show a wider margin
        \item The outlier point should be on the wrong side or inside the margin
        \item Draw a ``slack'' indicator from the outlier to where it ``should'' be:
        \begin{itemize}
            \item Dashed line from outlier perpendicular to margin boundary
            \item Label as $\xi_i$ (slack variable)
        \end{itemize}
        \item Shade the wider margin region
        \item Circle the support vectors (different from left panel)
        \item Add annotation: ``Wider margin, allows violations''
    \end{itemize}

    \item[Required Labels]
    \begin{itemize}
        \item Panel titles: ``Hard Margin SVM'' (left), ``Soft Margin SVM'' (right)
        \item Both panels: axes $x_1$, $x_2$
        \item Left panel: ``Outlier'', ``Narrow margin''
        \item Right panel: ``Slack $\xi_i$'', ``Wider margin''
        \item Legend (shared): Class +1 (blue), Class -1 (orange)
    \end{itemize}

    \item[Caption Text]
    ``Hard margin vs soft margin SVM. Left: Hard margin SVM requires perfect separation, so a single outlier (circled) forces a narrow, suboptimal margin. Right: Soft margin SVM allows margin violations controlled by slack variables $\xi_i$, achieving a wider margin at the cost of misclassifying the outlier. The trade-off is controlled by the regularization parameter $C$.''

    \item[Alt-Text for Accessibility]
    ``Two scatter plots comparing SVM approaches. Left panel shows hard margin SVM where an outlier point forces a tilted decision boundary with a narrow margin. Right panel shows soft margin SVM on the same data with a more centered decision boundary and wider margin, with the outlier allowed to violate the margin. A dashed line shows the slack distance from the outlier to the margin boundary.''
\end{description}

\newpage

%------------------------------------------------------------------------------
\subsection{Figure 9.7: Lagrangian Duality Visualization}
%------------------------------------------------------------------------------

\begin{description}[style=nextline, leftmargin=0pt]
    \item[Figure ID] \texttt{fig:lagrangian-duality}
    \item[Title] Lagrangian Duality: Primal Problem and Constraint
    \item[Type] 2D contour plot with constraint visualization
    \item[Size] 5.0 $\times$ 4.5 inches

    \item[Visual Description]
    \begin{itemize}
        \item Draw contour lines for a convex objective function (e.g., $f(x_1, x_2) = x_1^2 + x_2^2$)
        \item Domain: $x_1, x_2 \in [-2, 3]$
        \item Contour lines: concentric circles/ellipses, Primary Blue, various intensities
        \item Draw a linear constraint as a line (or curved constraint as a curve):
        \begin{itemize}
            \item Example: $g(x) = x_1 + x_2 - 1 \leq 0$
            \item Draw the constraint boundary as a thick line: Danger Red, 2.5pt
            \item Shade the feasible region (where constraint is satisfied) in Light Blue, 30\% opacity
            \item Shade the infeasible region with diagonal hatching or light red
        \end{itemize}
        \item Mark key points:
        \begin{itemize}
            \item Unconstrained minimum: center of contours, marked with $\times$ (gray), labeled ``Unconstrained optimum''
            \item Constrained minimum: point where lowest contour touches the constraint boundary
            \item Mark constrained minimum with a star (Success Green)
        \end{itemize}
        \item Draw the gradient vectors at the constrained minimum:
        \begin{itemize}
            \item $\nabla f$: pointing toward unconstrained minimum (Primary Blue arrow)
            \item $\nabla g$: pointing away from feasible region, perpendicular to constraint (Danger Red arrow)
            \item Show that they are parallel/anti-parallel (KKT condition: $\nabla f = \lambda \nabla g$)
        \end{itemize}
        \item Add annotation showing complementary slackness: ``$\lambda^* g(\mathbf{x}^*) = 0$''
    \end{itemize}

    \item[Required Labels]
    \begin{itemize}
        \item Axes: $x_1$, $x_2$
        \item Contour annotation: ``Level sets of $f(\mathbf{x})$''
        \item Constraint line: ``$g(\mathbf{x}) = 0$''
        \item Feasible region: ``Feasible region: $g(\mathbf{x}) \leq 0$''
        \item Unconstrained minimum: ``Unconstrained optimum $\mathbf{x}^*_{\text{unc}}$''
        \item Constrained minimum: ``Constrained optimum $\mathbf{x}^*$''
        \item Gradient arrows: $\nabla f(\mathbf{x}^*)$, $\lambda^* \nabla g(\mathbf{x}^*)$
    \end{itemize}

    \item[Caption Text]
    ``Geometric interpretation of Lagrangian duality. The contour lines show level sets of the objective function $f(\mathbf{x})$, with the unconstrained minimum at the center. The red line represents the constraint boundary $g(\mathbf{x}) = 0$, with the feasible region shaded. At the constrained optimum (green star), the gradients $\nabla f$ and $\nabla g$ are parallel---the KKT condition $\nabla f = \lambda \nabla g$ is satisfied.''

    \item[Alt-Text for Accessibility]
    ``2D contour plot illustrating constrained optimization. Concentric blue circular contour lines represent the objective function, centered at the unconstrained minimum marked with an X. A red diagonal line shows the constraint boundary, with the lower-left region shaded as feasible. The constrained optimum is marked with a green star where a contour line just touches the constraint. Two arrows at this point show the objective gradient and constraint gradient pointing in opposite directions.''
\end{description}

\newpage

%------------------------------------------------------------------------------
\subsection{Figure 9.8: Loss Function Comparison}
%------------------------------------------------------------------------------

\begin{description}[style=nextline, leftmargin=0pt]
    \item[Figure ID] \texttt{fig:loss-function-comparison}
    \item[Title] Classification Loss Functions: 0-1, Hinge, and Logistic Loss
    \item[Type] 2D line plot (function curves)
    \item[Size] 5.5 $\times$ 4.0 inches

    \item[Visual Description]
    \begin{itemize}
        \item Single plot with three loss functions plotted against the margin $m = y \cdot f(x)$
        \item X-axis (margin $m$): range $[-2, 3]$
        \item Y-axis (loss): range $[0, 3]$
        \item Grid: light gray, subtle
        \item Plot the following functions:

        \textbf{1. Zero-One Loss} (black, thick dashed):
        \begin{itemize}
            \item $L_{0-1}(m) = \mathbb{1}[m < 0] = \begin{cases} 1 & m < 0 \\ 0 & m \geq 0 \end{cases}$
            \item Step function: 1 for $m < 0$, 0 for $m \geq 0$
            \item Show discontinuity at $m = 0$ with open/closed circles
            \item Line style: thick dashed black, 2pt
        \end{itemize}

        \textbf{2. Hinge Loss} (Primary Blue, solid):
        \begin{itemize}
            \item $L_{\text{hinge}}(m) = \max(0, 1 - m)$
            \item Piecewise linear: decreasing line from $(−2, 3)$ to $(1, 0)$, then flat at 0
            \item Line style: solid, Primary Blue, 2.5pt
            \item Mark the ``corner'' at $m = 1$
        \end{itemize}

        \textbf{3. Logistic Loss} (Secondary Orange, solid):
        \begin{itemize}
            \item $L_{\text{log}}(m) = \log(1 + e^{-m})$
            \item Smooth curve, asymptotic to 0 as $m \to \infty$
            \item Line style: solid, Secondary Orange, 2.5pt
        \end{itemize}

        \item Add vertical dashed line at $m = 0$ (decision boundary)
        \item Add vertical dashed line at $m = 1$ (margin = 1 for SVM)
        \item Shade the region $m < 0$ lightly (incorrect classification zone)
    \end{itemize}

    \item[Required Labels]
    \begin{itemize}
        \item X-axis: ``Margin $m = y \cdot f(\mathbf{x})$''
        \item Y-axis: ``Loss $L(m)$''
        \item Legend (in plot area or beside):
        \begin{itemize}
            \item ``0-1 Loss'' (black dashed)
            \item ``Hinge Loss'' (blue solid)
            \item ``Logistic Loss'' (orange solid)
        \end{itemize}
        \item Vertical line at $m=0$: ``Decision boundary''
        \item Vertical line at $m=1$: ``Margin = 1''
        \item Region $m < 0$: ``Misclassified''
        \item Region $m > 0$: ``Correct''
    \end{itemize}

    \item[Caption Text]
    ``Comparison of classification loss functions. The 0-1 loss (dashed) is the true classification error but is non-convex and non-differentiable. The hinge loss (blue) used in SVMs is a convex upper bound that penalizes points inside the margin. The logistic loss (orange) used in logistic regression is smooth and convex. Both surrogate losses are convex relaxations of the intractable 0-1 loss.''

    \item[Alt-Text for Accessibility]
    ``Line plot comparing three loss functions against margin. The black dashed 0-1 loss is a step function, jumping from 1 to 0 at margin equals zero. The blue hinge loss decreases linearly from the left, reaching zero at margin equals 1 and staying at zero beyond. The orange logistic loss is a smooth curve that approaches zero asymptotically for positive margins. Vertical dashed lines mark the decision boundary at margin zero and the SVM margin at one.''
\end{description}

\newpage

%==============================================================================
\section{Chapter 15: Convergence Analysis (6 Figures)}
%==============================================================================

%------------------------------------------------------------------------------
\subsection{Figure 15.1: Convergence Rate Comparison}
%------------------------------------------------------------------------------

\begin{description}[style=nextline, leftmargin=0pt]
    \item[Figure ID] \texttt{fig:convergence-rate-comparison}
    \item[Title] Convergence Rate Comparison: Sublinear vs Linear
    \item[Type] Semi-log plot (log scale on y-axis)
    \item[Size] 5.5 $\times$ 4.0 inches

    \item[Visual Description]
    \begin{itemize}
        \item Plot showing convergence of different rates on a log scale
        \item X-axis (iterations $t$): linear scale, range $[0, 200]$
        \item Y-axis (suboptimality $f(\theta_t) - f(\theta^*)$): logarithmic scale, range $[10^{-8}, 10^0]$
        \item Grid: horizontal lines at each power of 10, vertical lines every 50 iterations
        \item Plot the following convergence curves:

        \textbf{1. Sublinear $O(1/t)$} (Danger Red):
        \begin{itemize}
            \item Formula: $1/(t+1)$
            \item Line style: solid, 2pt
            \item This should be the slowest-converging curve
            \item On log scale, appears as gradually curving downward
        \end{itemize}

        \textbf{2. Accelerated Sublinear $O(1/t^2)$} (Secondary Orange):
        \begin{itemize}
            \item Formula: $1/(t+1)^2$
            \item Line style: solid, 2pt
            \item Faster than $O(1/t)$, steeper curve
        \end{itemize}

        \textbf{3. Linear $O(\rho^t)$ with $\rho = 0.95$} (Primary Blue):
        \begin{itemize}
            \item Formula: $0.95^t$
            \item Line style: solid, 2.5pt
            \item On log scale, appears as a straight line
            \item Much faster convergence
        \end{itemize}

        \textbf{4. Linear $O(\rho^t)$ with $\rho = 0.9$} (Success Green):
        \begin{itemize}
            \item Formula: $0.9^t$
            \item Line style: dashed, 2pt
            \item Steeper straight line (faster convergence)
        \end{itemize}

        \item Add a horizontal dashed line at $\epsilon = 10^{-6}$ (target accuracy)
        \item Mark where each curve crosses this threshold
    \end{itemize}

    \item[Required Labels]
    \begin{itemize}
        \item X-axis: ``Iterations $t$''
        \item Y-axis: ``Suboptimality $f(\boldsymbol{\theta}_t) - f(\boldsymbol{\theta}^*)$''
        \item Legend:
        \begin{itemize}
            \item ``$O(1/t)$ -- GD on convex'' (red)
            \item ``$O(1/t^2)$ -- Accelerated GD'' (orange)
            \item ``$O(0.95^t)$ -- GD on strongly convex'' (blue)
            \item ``$O(0.9^t)$ -- Better conditioned'' (green)
        \end{itemize}
        \item Horizontal line: ``Target accuracy $\epsilon$''
        \item Annotations at crossings: iteration counts to reach $\epsilon$
    \end{itemize}

    \item[Caption Text]
    ``Comparison of convergence rates on a semi-log plot. Linear convergence (blue, green) appears as a straight line on the log scale and reaches machine precision in tens to hundreds of iterations. Sublinear convergence (red, orange) curves and requires orders of magnitude more iterations. The difference between $O(1/t)$ and $O(\rho^t)$ is the difference between minutes and hours of training.''

    \item[Alt-Text for Accessibility]
    ``Semi-log plot comparing four convergence curves over 200 iterations. Two curves showing linear convergence (blue and green) appear as straight downward lines, reaching very small values quickly. Two curves showing sublinear convergence (red and orange) appear as gradually curving lines that descend much more slowly. A horizontal dashed line marks the target accuracy level.''
\end{description}

\newpage

%------------------------------------------------------------------------------
\subsection{Figure 15.2: Condition Number Effect on Convergence}
%------------------------------------------------------------------------------

\begin{description}[style=nextline, leftmargin=0pt]
    \item[Figure ID] \texttt{fig:condition-number-effect}
    \item[Title] Effect of Condition Number on Optimization
    \item[Type] Side-by-side 2D contour plots with trajectories
    \item[Size] 6.0 $\times$ 3.0 inches (two panels)

    \item[Visual Description]
    Two panels showing gradient descent on quadratic functions with different condition numbers:

    \textbf{Left Panel -- ``Well-Conditioned ($\kappa = 1$)'':}
    \begin{itemize}
        \item Plot contours of $f(\theta_1, \theta_2) = \theta_1^2 + \theta_2^2$
        \item Contours are perfect circles (equal scaling in both directions)
        \item Domain: $\theta_1, \theta_2 \in [-2, 2]$
        \item Contour colors: gradient from dark (center) to light
        \item Show gradient descent trajectory:
        \begin{itemize}
            \item Start from $(1.5, 1.5)$
            \item Path goes almost directly to the center
            \item Only 5--8 iterations needed
            \item Line: Primary Blue, 2pt, with arrows
        \end{itemize}
        \item Mark start (red circle) and end (green star)
        \item Annotation: ``Direct path to minimum''
    \end{itemize}

    \textbf{Right Panel -- ``Ill-Conditioned ($\kappa = 100$)'':}
    \begin{itemize}
        \item Plot contours of $f(\theta_1, \theta_2) = \theta_1^2 + 100\theta_2^2$
        \item Contours are highly elongated ellipses (100:1 aspect ratio)
        \item Same domain as left panel
        \item Show gradient descent trajectory:
        \begin{itemize}
            \item Start from same relative position
            \item Path shows pronounced zigzag/oscillation pattern
            \item Many more iterations needed (20--30 shown)
            \item The trajectory bounces back and forth across the narrow valley
            \item Line: Primary Blue, 2pt, with arrows
        \end{itemize}
        \item Mark start (red circle) and end (green star)
        \item Annotation: ``Zigzag oscillations slow convergence''
    \end{itemize}

    \item[Required Labels]
    \begin{itemize}
        \item Panel titles: ``$\kappa = L/\mu = 1$'' (left), ``$\kappa = L/\mu = 100$'' (right)
        \item Both panels: axes $\theta_1$, $\theta_2$
        \item Start point: ``Start''
        \item End point: ``$\boldsymbol{\theta}^*$''
        \item Left: ``$\sim$8 iterations''
        \item Right: ``$\sim$500+ iterations''
    \end{itemize}

    \item[Caption Text]
    ``Effect of condition number $\kappa = L/\mu$ on gradient descent convergence. Left: When $\kappa = 1$ (circular contours), gradient descent takes a nearly direct path to the minimum. Right: When $\kappa = 100$ (elongated ellipses), gradient descent zigzags across the narrow valley, requiring $O(\kappa)$ times more iterations. This explains why ill-conditioned problems train slowly.''

    \item[Alt-Text for Accessibility]
    ``Two contour plots comparing gradient descent behavior. Left panel shows circular contours with a gradient descent path going nearly straight from the starting point to the center in about 8 steps. Right panel shows highly elongated elliptical contours with a gradient descent path that zigzags back and forth many times while slowly approaching the center, illustrating the difficulty of optimizing ill-conditioned problems.''
\end{description}

\newpage

%------------------------------------------------------------------------------
\subsection{Figure 15.3: Smoothness Illustration}
%------------------------------------------------------------------------------

\begin{description}[style=nextline, leftmargin=0pt]
    \item[Figure ID] \texttt{fig:smoothness-illustration}
    \item[Title] Smooth vs Non-Smooth Functions
    \item[Type] Side-by-side 1D function plots
    \item[Size] 6.0 $\times$ 3.0 inches (two panels)

    \item[Visual Description]
    Two panels comparing smooth and non-smooth functions:

    \textbf{Left Panel -- ``$L$-Smooth Function'':}
    \begin{itemize}
        \item Plot a smooth convex function: $f(x) = x^2$ or $f(x) = \frac{1}{2}x^2$
        \item Domain: $x \in [-2, 2]$
        \item Function curve: Primary Blue, 2.5pt solid
        \item At a point $x_0$ (e.g., $x_0 = 1$):
        \begin{itemize}
            \item Draw the tangent line (green dashed)
            \item Draw the quadratic upper bound: $f(x_0) + f'(x_0)(x - x_0) + \frac{L}{2}(x-x_0)^2$
            \item This is a parabola that touches $f$ at $x_0$ and curves upward
            \item Color: Secondary Orange, dashed
        \end{itemize}
        \item The function should lie between the tangent line (below) and quadratic bound (above)
        \item Mark $x_0$ with a red dot
        \item Shade the region between tangent and quadratic bound lightly
        \item Annotation: ``Gradient changes smoothly''
    \end{itemize}

    \textbf{Right Panel -- ``Non-Smooth Function'':}
    \begin{itemize}
        \item Plot $f(x) = |x|$ (absolute value)
        \item Domain: $x \in [-2, 2]$
        \item Function curve: Primary Blue, 2.5pt solid
        \item Show the ``kink'' at $x = 0$ clearly
        \item Draw two different tangent slopes at $x = 0$:
        \begin{itemize}
            \item Left-side slope: $-1$ (green dashed line extending left)
            \item Right-side slope: $+1$ (green dashed line extending right)
        \end{itemize}
        \item Mark the non-differentiable point at origin with a red circle
        \item Show subgradient interval: small vertical bracket at $x=0$ from $-1$ to $+1$
        \item Annotation: ``Gradient undefined at kink''
    \end{itemize}

    \item[Required Labels]
    \begin{itemize}
        \item Panel titles: ``Smooth: $f(x) = x^2$'' (left), ``Non-smooth: $f(x) = |x|$'' (right)
        \item Both panels: axes $x$, $f(x)$
        \item Left panel:
        \begin{itemize}
            \item ``Tangent line''
            \item ``Quadratic upper bound''
            \item ``$x_0$'' at tangent point
        \end{itemize}
        \item Right panel:
        \begin{itemize}
            \item ``Subgradient $\partial f(0) = [-1, 1]$''
            \item ``Non-differentiable point''
        \end{itemize}
    \end{itemize}

    \item[Caption Text]
    ``Smooth vs non-smooth functions. Left: An $L$-smooth function has a gradient that changes continuously. At any point, the function is sandwiched between its tangent line (lower bound) and a quadratic upper bound. This enables gradient descent with predictable step sizes. Right: The absolute value function $|x|$ has a kink at the origin where the gradient is undefined. At non-smooth points, we use subgradients instead, which complicates optimization.''

    \item[Alt-Text for Accessibility]
    ``Two function plots. Left panel shows a smooth parabola with a tangent line touching the curve at one point and a dashed quadratic upper bound curve above it, illustrating the smoothness property. Right panel shows the absolute value function forming a V-shape with two different tangent lines meeting at the sharp corner at the origin, illustrating a non-smooth point where the gradient is undefined.''
\end{description}

\newpage

%------------------------------------------------------------------------------
\subsection{Figure 15.4: Momentum Effect on Optimization}
%------------------------------------------------------------------------------

\begin{description}[style=nextline, leftmargin=0pt]
    \item[Figure ID] \texttt{fig:momentum-effect}
    \item[Title] Gradient Descent vs Momentum: Trajectory Comparison
    \item[Type] 2D contour plot with two trajectories
    \item[Size] 5.5 $\times$ 4.5 inches

    \item[Visual Description]
    \begin{itemize}
        \item Draw elliptical contour lines for an ill-conditioned quadratic ($\kappa \approx 10$)
        \item Function: $f(\theta_1, \theta_2) = \theta_1^2 + 10\theta_2^2$
        \item Domain: $\theta_1, \theta_2 \in [-2.5, 2.5]$
        \item Contour lines: Neutral Gray with gradient fill
        \item Mark the global minimum at origin with a green star

        \item \textbf{Trajectory 1 -- Vanilla Gradient Descent} (Danger Red):
        \begin{itemize}
            \item Starting point: $(2, 1)$
            \item Show characteristic zigzag pattern
            \item Many iterations (15--20 visible)
            \item Line: Danger Red, 2pt solid, with small circles at each iteration
            \item The zigzag should bounce perpendicular to the elongated contours
        \end{itemize}

        \item \textbf{Trajectory 2 -- Gradient Descent with Momentum} (Primary Blue):
        \begin{itemize}
            \item Same starting point: $(2, 1)$
            \item Show smooth, accelerated path
            \item Fewer iterations needed (8--10 visible)
            \item Line: Primary Blue, 2.5pt solid, with small circles at each iteration
            \item Path should show initial acceleration then smooth approach
            \item May slightly overshoot then correct
        \end{itemize}

        \item Both trajectories should start from the same point but take different paths
        \item Mark the starting point prominently (filled black circle)
    \end{itemize}

    \item[Required Labels]
    \begin{itemize}
        \item Axes: $\theta_1$, $\theta_2$
        \item Starting point: ``Start''
        \item Minimum: ``$\boldsymbol{\theta}^*$''
        \item Legend:
        \begin{itemize}
            \item ``Vanilla GD'' (red) -- ``20 iterations''
            \item ``GD + Momentum'' (blue) -- ``8 iterations''
        \end{itemize}
        \item Annotation on red path: ``Zigzag oscillations''
        \item Annotation on blue path: ``Smooth acceleration''
    \end{itemize}

    \item[Caption Text]
    ``Momentum eliminates oscillations on ill-conditioned problems. Vanilla gradient descent (red) zigzags across the narrow valley, wasting iterations on perpendicular oscillations. Gradient descent with momentum (blue) accumulates velocity along the consistent gradient direction, damping the oscillations and converging in fewer iterations. Momentum is especially beneficial when $\kappa$ is large.''

    \item[Alt-Text for Accessibility]
    ``Contour plot showing elongated elliptical level curves with two optimization paths starting from the same point. The red path labeled vanilla gradient descent shows a pronounced zigzag pattern, oscillating many times before reaching the minimum. The blue path labeled gradient descent with momentum shows a much smoother curved trajectory that reaches the minimum directly with fewer iterations.''
\end{description}

\newpage

%------------------------------------------------------------------------------
\subsection{Figure 15.5: SGD vs GD Trajectories}
%------------------------------------------------------------------------------

\begin{description}[style=nextline, leftmargin=0pt]
    \item[Figure ID] \texttt{fig:sgd-vs-gd}
    \item[Title] Stochastic Gradient Descent vs Full-Batch Gradient Descent
    \item[Type] 2D contour plot with two trajectories
    \item[Size] 5.5 $\times$ 4.5 inches

    \item[Visual Description]
    \begin{itemize}
        \item Draw contour lines for a moderately ill-conditioned function
        \item Domain: $\theta_1, \theta_2 \in [-2, 2]$
        \item Contour lines: light gray with subtle gradient fill
        \item Mark the global minimum with a green star

        \item \textbf{Trajectory 1 -- Full-Batch GD} (Primary Blue):
        \begin{itemize}
            \item Starting point: $(1.5, 1.5)$
            \item Smooth, deterministic path
            \item Clean steps perpendicular to contours
            \item Line: Primary Blue, 2.5pt solid
            \item Show 10--12 iterations with clear markers
            \item Path converges precisely to the minimum
        \end{itemize}

        \item \textbf{Trajectory 2 -- Stochastic GD} (Secondary Orange):
        \begin{itemize}
            \item Same starting point
            \item Noisy, jittery path
            \item General direction is correct but with random perturbations
            \item Line: Secondary Orange, 2pt, slightly more iterations shown
            \item Each step should deviate randomly from the true gradient direction
            \item Path converges to a neighborhood around the minimum, not exactly to it
            \item Show the path ``wandering'' near the end rather than settling precisely
        \end{itemize}

        \item Add a small shaded circle around the minimum showing SGD's convergence neighborhood
        \item Mark the starting point with a black filled circle
    \end{itemize}

    \item[Required Labels]
    \begin{itemize}
        \item Axes: $\theta_1$, $\theta_2$
        \item Starting point: ``Start $\boldsymbol{\theta}_0$''
        \item Minimum: ``$\boldsymbol{\theta}^*$''
        \item Legend:
        \begin{itemize}
            \item ``Full-batch GD'' (blue)
            \item ``Stochastic GD'' (orange)
        \end{itemize}
        \item Shaded region: ``SGD convergence neighborhood''
        \item Annotation on GD path: ``Smooth, deterministic''
        \item Annotation on SGD path: ``Noisy, stochastic''
    \end{itemize}

    \item[Caption Text]
    ``SGD vs full-batch gradient descent trajectories. Full-batch GD (blue) follows a smooth path using the exact gradient computed over all data. SGD (orange) uses noisy gradient estimates from mini-batches, resulting in a jittery path. SGD converges to a neighborhood around the optimum whose size depends on the learning rate and batch size, not to the exact minimum.''

    \item[Alt-Text for Accessibility]
    ``Contour plot with two optimization trajectories. The blue full-batch gradient descent path follows a smooth curved line from the starting point directly to the minimum at the center. The orange stochastic gradient descent path follows a noisy, jagged trajectory that generally moves toward the minimum but wanders randomly around the final destination. A small shaded circle around the minimum indicates SGD's convergence region.''
\end{description}

\newpage

%------------------------------------------------------------------------------
\subsection{Figure 15.6: Learning Rate Failure Modes}
%------------------------------------------------------------------------------

\begin{description}[style=nextline, leftmargin=0pt]
    \item[Figure ID] \texttt{fig:learning-rate-failures}
    \item[Title] Learning Rate Failure Modes
    \item[Type] Loss curves over training iterations
    \item[Size] 5.5 $\times$ 4.0 inches

    \item[Visual Description]
    \begin{itemize}
        \item Single plot showing loss vs iteration for different learning rates
        \item X-axis (iterations): range $[0, 100]$
        \item Y-axis (loss): range $[0, 5]$ (adjust as needed)
        \item Light grid lines

        \item \textbf{Curve 1 -- Learning Rate Too High} (Danger Red):
        \begin{itemize}
            \item Start at moderate loss value
            \item Show divergent/unstable behavior
            \item Loss oscillates wildly or increases over time
            \item Eventually goes to infinity or explodes
            \item Line: Danger Red, 2pt solid
            \item May need to truncate at top of plot with arrow indicating ``$\to \infty$''
        \end{itemize}

        \item \textbf{Curve 2 -- Learning Rate Slightly Too High} (Secondary Orange):
        \begin{itemize}
            \item Start at same loss value
            \item Show oscillating convergence (damped oscillations)
            \item Loss bounces around but eventually decreases
            \item Convergence is slow and inefficient
            \item Line: Secondary Orange, 2pt solid
        \end{itemize}

        \item \textbf{Curve 3 -- Optimal Learning Rate} (Success Green):
        \begin{itemize}
            \item Start at same loss value
            \item Show smooth, fast convergence
            \item Monotonically decreasing (or nearly so)
            \item Reaches low loss value quickly and stabilizes
            \item Line: Success Green, 2.5pt solid
        \end{itemize}

        \item \textbf{Curve 4 -- Learning Rate Too Low} (Primary Blue):
        \begin{itemize}
            \item Start at same loss value
            \item Show very slow, gradual decrease
            \item Smooth but painfully slow progress
            \item After 100 iterations, still hasn't reached the good loss value
            \item Line: Primary Blue, 2pt solid
        \end{itemize}

        \item Add horizontal dashed line at the target/optimal loss value
    \end{itemize}

    \item[Required Labels]
    \begin{itemize}
        \item X-axis: ``Iteration''
        \item Y-axis: ``Loss $\mathcal{L}(\boldsymbol{\theta})$''
        \item Legend:
        \begin{itemize}
            \item ``$\eta$ too high -- diverges'' (red)
            \item ``$\eta$ slightly high -- oscillates'' (orange)
            \item ``$\eta$ optimal -- fast convergence'' (green)
            \item ``$\eta$ too low -- slow convergence'' (blue)
        \end{itemize}
        \item Horizontal line: ``Optimal loss''
        \item Annotation on red curve: ``Divergence''
        \item Annotation on blue curve: ``Slow progress''
    \end{itemize}

    \item[Caption Text]
    ``Learning rate failure modes. When the learning rate $\eta$ is too high (red), optimization diverges as updates overshoot the minimum. A slightly high rate (orange) causes oscillations but may eventually converge. The optimal rate (green) achieves fast, stable convergence. A rate that is too low (blue) converges smoothly but wastes computation. The optimal learning rate is $\eta = O(1/L)$ where $L$ is the smoothness constant.''

    \item[Alt-Text for Accessibility]
    ``Line plot showing four loss curves over 100 iterations. The red curve (learning rate too high) oscillates wildly and trends upward, indicating divergence. The orange curve (slightly high) oscillates but eventually decreases. The green curve (optimal) decreases smoothly and quickly to a low value. The blue curve (too low) decreases very slowly and remains above the optimal value even after 100 iterations. A horizontal dashed line marks the optimal loss level.''
\end{description}

\newpage

%==============================================================================
\section{Production Checklist}
%==============================================================================

\begin{enumerate}
    \item \textbf{Color consistency}: Verify all figures use the defined color palette
    \item \textbf{Font consistency}: Use same font family across all figures
    \item \textbf{Line weights}: Primary elements 2--2.5pt, secondary elements 1.5pt, grid 0.5pt
    \item \textbf{Mathematical notation}: Use proper LaTeX-rendered symbols
    \item \textbf{Resolution}: Export at 300 DPI minimum for print
    \item \textbf{File naming}: Use figure IDs as filenames (e.g., \texttt{fig-convex-nonconvex-sets.pdf})
    \item \textbf{Formats}: Provide both PDF (vector) and PNG (300 DPI) versions
    \item \textbf{Alt-text}: Include accessibility descriptions in figure metadata
    \item \textbf{Bleed}: Ensure no critical elements within 0.125'' of edges
    \item \textbf{Color mode}: Use CMYK for print, RGB for digital versions
\end{enumerate}

%==============================================================================
\section{Summary Table}
%==============================================================================

\begin{longtable}{p{1.5cm}p{4cm}p{3cm}p{5cm}}
\toprule
\textbf{Figure} & \textbf{Title} & \textbf{Type} & \textbf{Key Elements} \\
\midrule
\endhead
\multicolumn{4}{l}{\textbf{Chapter 9: Convex Optimization}} \\
\midrule
9.1 & Convex vs Non-Convex Sets & 2D comparison & Line segment test, two shapes \\
9.2 & Convex vs Concave Functions & 3D surfaces & Bowl and dome, chords \\
9.3 & First-Order Condition & 3D surface + plane & Tangent plane as underestimator \\
9.4 & Gradient Descent on Convex & 2D contours & Trajectory, gradient arrows \\
9.5 & SVM Decision Boundary & 2D scatter & Margin, support vectors \\
9.6 & Hard vs Soft Margin SVM & 2D comparison & Outlier handling, slack variables \\
9.7 & Lagrangian Duality & 2D contours & Constraint, KKT gradients \\
9.8 & Loss Function Comparison & 1D curves & 0-1, hinge, logistic losses \\
\midrule
\multicolumn{4}{l}{\textbf{Chapter 15: Convergence Analysis}} \\
\midrule
15.1 & Convergence Rate Comparison & Semi-log plot & $O(1/t)$ vs $O(\rho^t)$ curves \\
15.2 & Condition Number Effect & 2D comparison & $\kappa=1$ vs $\kappa=100$ contours \\
15.3 & Smoothness Illustration & 1D comparison & Smooth vs non-smooth functions \\
15.4 & Momentum Effect & 2D contours & GD zigzag vs momentum smooth \\
15.5 & SGD vs GD Trajectories & 2D contours & Clean vs noisy paths \\
15.6 & Learning Rate Failure Modes & Loss curves & Diverge, oscillate, optimal, slow \\
\bottomrule
\end{longtable}

\end{document}
